% Páginas de entrada comunes a los libros de Comprensión lectora (3 y 4)
\cleardoublepage
\thispagestyle{vacia}
{\sffamily\fontsize{26}{30}\selectfont\firalight Cómo se usa este libro\par}\vspace{2pt}
{\color{zolive}\rule{18mm}{1.6pt}}\par\vspace{8pt}

Leer bien no es leer rápido ni leer mucho: es saber \emph{qué preguntarle} a un texto. Este libro enseña, una por una, las
operaciones que la prueba de admisión evalúa: encontrar un dato, reconocer una paráfrasis, deducir lo que no se dice,
entender una palabra por su contexto, seguir una referencia, identificar la tesis de un autor. Cada lección del libro
corresponde a una lección de Comprensión lectora en la app \marca, con el mismo código.

\vspace{4pt}
\begin{tabularx}{\linewidth}{@{}>{\leavevmode\sffamily\small\bfseries\color{zolive}}p{27mm}X@{}}
La ruta & Abre la lección: los pasos que vas a recorrer, los mismos de la app. El último es siempre un problema.\\[4pt]
La analogía & Una comparación con algo cotidiano para entender la técnica antes de aplicarla.\\[4pt]
Cuadro por cuadro & Cada idea se explica sobre un texto marcado: a la izquierda, la hoja con lo que se subraya, se tacha o se
anota al margen; a la derecha, los pasos numerados.\\[4pt]
Tu turno & Enseguida, un ejercicio con la misma hoja para marcar con lápiz. Las hojas resueltas están en \emph{Soluciones
dibujadas}.\\[4pt]
Ejemplo resuelto & Una pregunta como las del examen, con el texto marcado y cada opción analizada.\\[4pt]
La técnica y la regla & El procedimiento con todas las letras y el recuadro para llevar al examen.\\[4pt]
El texto & Una lectura original, escrita para este libro, con los párrafos numerados (\S\,1, \S\,2…), como en la
prueba. Se lee con lápiz: subraya la idea de cada párrafo.\\[4pt]
Tu mapa del texto & Después de cada lectura, una palabra por párrafo (presenta, explica, objeta…) y la idea central en una
oración. Es lo que harás en el examen en un minuto.\\[4pt]
Preguntas comentadas & Primero ves la pregunta y las opciones y escribes tu respuesta. Debajo de la línea de tijera está el
análisis de cada opción: por qué la correcta lo es y por qué cada distractor falla. Aquí se aprende a \emph{descartar}.\\[4pt]
Ponlo en práctica & Con la técnica ya entendida, un reto tipo examen para resolver tú solo. Debajo de la línea de tijera,
«Así se usó el concepto» muestra cómo se aplica la técnica paso a paso.\\[4pt]
Ahora tú & Fragmentos nuevos para aplicar la técnica tú solo, con renglones para responder. Las respuestas están en las
claves.\\[4pt]
Trampa y atajo & El distractor favorito de ese tipo de pregunta y una manera rápida de evitarlo.\\[4pt]
Taller & Preguntas abiertas para escribir y preguntas de opción múltiple como las del examen. Las respuestas, con su
justificación, están al final, en \emph{Claves y soluciones}.\\[4pt]
Evaluación & Cada mundo cierra con un texto nuevo y sus preguntas, con tiempo.\\
\end{tabularx}

\vspace{6pt}
\begin{tcolorbox}[enhanced,colback=zsoft,coltext=zink,boxrule=0pt,arc=3pt,left=5mm,right=5mm,top=3mm,bottom=3mm]
{\sffamily\scriptsize\bfseries\color{zolive}CUATRO TIPOS DE DISTRACTOR}\par\vspace{3pt}
\small Casi todas las opciones falsas del examen caen en una de estas familias. Aprende a nombrarlas y las verás venir.
\begin{itemize}
  \item[\color{zolive}1] \textbf{Dice de más:} exagera o generaliza lo que el texto dice con matices («siempre», «todos»).
  \item[\color{zolive}2] \textbf{Dice otra cosa:} usa palabras del texto pero cambia la relación (causa por efecto, el
  autor por otra voz).
  \item[\color{zolive}3] \textbf{No lo dice:} puede ser cierto en el mundo, pero el texto no lo afirma ni lo permite
  deducir.
  \item[\color{zolive}4] \textbf{Dice lo contrario:} invierte el sentido, a veces con una negación escondida.
\end{itemize}
\end{tcolorbox}

\clearpage
\thispagestyle{vacia}
{\sffamily\fontsize{26}{30}\selectfont\firalight La prueba de lectura\par}\vspace{2pt}
{\color{zolive}\rule{18mm}{1.6pt}}\par\vspace{8pt}

El componente de Competencia lectora de la prueba de admisión de la Universidad de Antioquia vale lo mismo que el de
Razonamiento lógico: alrededor de \textbf{40 preguntas} de opción múltiple sobre textos de varios párrafos (ensayos,
artículos de divulgación, crónicas, relatos, cartas, a veces dos textos que dialogan). Las preguntas recorren tres niveles:

\begin{enumerate}[leftmargin=1.6em,itemsep=4pt]
  \item \textbf{Literal:} lo que el texto dice. Ubicar un dato, reconocer una paráfrasis, saber quién dice qué.
  \item \textbf{Inferencial:} lo que el texto da a entender. Ideas implícitas, causas, actitudes, el sentido de una palabra o
  de un conector, la función de un párrafo.
  \item \textbf{Crítico:} lo que el texto hace. La tesis y los argumentos, la intención y el tono, cómo se aplicaría su idea
  a otro caso, qué diría un autor del otro.
\end{enumerate}

Este Libro 3 se concentra en los dos primeros niveles («leer lo que dice»); el Libro 4, en el segundo y el tercero
(«leer lo que quiere decir»).

\subsection*{Un método para cada texto}
\begin{enumerate}
  \item \textbf{Prelee} (30 segundos): título, primera y última frase de cada párrafo, fuente.
  \item \textbf{Lee con lápiz}: una palabra al margen de cada párrafo que diga de qué trata (su «función»).
  \item \textbf{Lee la pregunta} antes que las opciones, y decide qué tipo de pregunta es.
  \item \textbf{Vuelve al texto}: casi ninguna respuesta se contesta de memoria. El párrafo exacto tiene la prueba.
  \item \textbf{Descarta} con las cuatro familias de distractores. Quédate con la opción que el texto \emph{sostiene}, no
  con la que te parece verdadera.
\end{enumerate}

\subsection*{El tiempo}
Con dos textos largos y unas veinte preguntas por texto, conviene dedicar a cada texto unos 6 minutos de lectura y luego
alrededor de un minuto por pregunta. Si una pregunta se atasca, márcala y sigue: muchas veces otra pregunta posterior te
obliga a releer justo el párrafo que necesitabas.
